Las catorce fases del proyecto y los requisitos principales que cubre
cada una se resumen en la \cref{table:impl-phases}. % TODO cuadro en vez de tabla

\begin{table}[Fases del proyecto]%
            {table:impl-phases}%
            {Plan de trabajo en catorce fases.}
  \small
  \renewcommand{\arraystretch}{1.3}
  \begin{tabular}{c p{7cm} p{4cm}}
    \hline
    \textbf{Fase} & \textbf{Nombre} & \textbf{Requisitos principales} \\
    \hline
    0  & Configuración del proyecto y arquitectura base
       & RNF-3, RNF-9, RNF-13 \\
    1  & Autenticación y organizaciones
       & RF-1, RF-2, RF-16.3 \\
    2  & Cursos académicos
       & RF-3 \\
    3  & Asignaturas, grupos y membresías
       & RF-4, RF-5 \\
    4  & Exámenes, modelos, perfiles de página y \acs{pdf} instrumentado
       & RF-6 \\
    5  & Instancias de examen, máquina de estados y corrección
       & RF-7, RF-8, RF-11 \\
    6  & Ingesta, dispatcher de reconocedores y ensamblado
       & RF-9 \\
    7  & Anotaciones (texto, stylus, audio)
       & RF-10 \\
    8  & Revisiones de examen
       & RF-12 \\
    9  & Notificaciones y correos electrónicos
       & RF-13 \\
    10 & Exportación, búsqueda avanzada y configuración
       & RF-14, RF-15 \\
    11 & Tareas de mantenimiento y borrado automático
       & RF-15.3, RF-15.4 \\
    12 & Seguridad avanzada, protección de páginas y pruebas
       & RF-16, RNF-6, RNF-12 \\
    13 & Preparación para \acs{sso} y despliegue final
       & RF-1.5, RNF-3 \\
    \hline
  \end{tabular}
\end{table}

La distribución refleja el tamaño relativo de cada fase, calculado a
partir de dos métricas: el número de líneas de código asociadas al
dominio funcional de la fase y el número de subrequisitos que cubre. Las
fases más voluminosas son la 1 (autenticación y organizaciones), la 5
(instancias y máquina de estados de corrección) y la 12 (seguridad
avanzada y pruebas, de carácter transversal).
